\SeventhArticle{19}
Da es genügt, eine der beiden Funktionen $T,U$ zu kennen, haben wir die Bestimmung der letzteren als die einfachere vorgezogen. Wie diese sich auf die Entwicklung der Reihe $u^{-1}+\frac13u^{-3}+\frac15u^{-5}+\text{usw.}$ in einen Kettenbruch stützt, hätten wir durch ähnliche Überlegungen aus der Entwicklung der Reihe $t^{-1}+\frac12t^{-2}+\frac13t^{-3}+\frac14t^{-4}+\text{usw.}$ in den Kettenbruch
\SourceLeaf{199}
\[
\cfrac1{t-\cfrac{\frac12}{1-\cfrac{\frac16}
{t-\cfrac{\frac26}{1-\cfrac{\frac2{10}}
{t-\cfrac{\frac3{10}}{1-\mathrm{usw.}}}}}}}
\]
einen Algorithmus zur Bestimmung der Funktion $T$ für die aufeinanderfolgenden Werte der Zahl $n$ herleiten können. Zu demselben Ergebnis gelangen wir jedoch, wenn wir bedenken, dass $T$ nichts anderes als $\frac U{2^{n+1}}$ oder $\frac{W^{(n+1)}}{2^{n+1}}$ ist, wenn für $u$ der Ausdruck $2t-1$ geschrieben wird. So erhält man die nacheinander für $T$ zu verwendenden Funktionen durch den folgenden Algorithmus:
\[
\begin{aligned}
W&=1\\
\tfrac12W'&=t-\tfrac12\\
\tfrac14W''&=(t-\tfrac12).\tfrac12W'-\frac{1.1}{2.6}W
\\[3pt]&\qquad=tt-t+\tfrac16\\
\tfrac18W'''&=(t-\tfrac12).\tfrac14W''-\frac{2.2}{6.10}.\tfrac12W'
\\[3pt]&\qquad=t^3-\tfrac32tt+\tfrac35t-\tfrac1{20}\\
\tfrac1{10}W''''&=(t-\tfrac12).\tfrac18W'''-\frac{3.3}{10.14}.\tfrac14W''
\\[3pt]&\qquad=t^4-2t^3+\tfrac97tt-\tfrac27t+\tfrac1{70}
\end{aligned}
\]
usw. Durch Induktion ergibt sich daraus allgemein
\[
\begin{aligned}
T={}&t^{n+1}-\frac{(n+1)^2}{1.(2n+2)}t^n
+\frac{(n+1)^2.nn}{1.2.(2n+2)(2n+1)}t^{n-1}\\
&-\frac{(n+1)^2.nn.(n-1)^2}{1.2.3.(2n+2)(2n+1).2n}t^{n-2}
+\mathrm{usw.}
\end{aligned}
\]
oder $T=t^{n+1}F(-(n+1),-(n+1),-2(n+1),t^{-1})$; dieser Induktion lässt sich leicht die Kraft eines Beweises verleihen. Wenn man es vorzieht, kann $T$ auch bei umgekehrter Reihenfolge der Glieder dargestellt werden durch
\[
\pm\frac{1.2.3.4\mathinner{.\,.\,.\,.\,.}(n+1)}
{2.6.10.14\mathinner{.\,.\,.\,.\,.}(4n+2)}
F(n+2,-(n+1),1,t)
\]
wobei das obere Zeichen für ungerades $n$, das untere für gerades gilt. Schließlich ergibt sich auf ähnliche Weise allgemein $T''$ als Produkt aus
\[
\frac{1.1.2.2.3.3\mathinner{.\,.\,.\,.\,.}(n+1).(n+1)}
{2.6.6.10.10.14\mathinner{.\,.\,.\,.\,.}(4n+2).(4n+6)}t^{-(n+2)}
\]
und der unendlichen Reihe
\[
\begin{aligned}
1&+\frac{(n+2)^2}{1.(2n+4)}t^{-1}
+\frac{(n+2)^2(n+3)^2}{1.2.(2n+4)(2n+5)}t^{-2}\\
&+\frac{(n+2)^2.(n+3)^2(n+4)^2}
{1.2.3.(2n+4)(2n+5)(2n+6)}t^{-3}+\mathrm{usw.}
\end{aligned}
\]
oder $F(n+2,n+2,2n+4,t^{-1})$
\SourceLeaf{200}
\SeventhArticle{20}
Da in der Funktion $U$ die Potenzen $u^n$, $u^{n-2}$, $u^{n-4}$ usw. fehlen, werden unter den Wurzeln der Gleichung $U=0$ stets je zwei von gleicher Größe und entgegengesetztem Vorzeichen sein; bei geradem Wert von $n$ muss außerdem die einzelne Wurzel $0$ hinzukommen. Sind die Wurzeln gefunden, so erhält man die Werte der Koeffizienten $R,R',R''$ usw. nach dem Verfahren des Artikels 11 durch ein Polynom in $u$, das bei ungeradem Wert von $n$ die Form
\[
\gamma u^{n-1}+\gamma'u^{n-3}+\gamma''u^{n-5}+\mathrm{usw.}
\]
und bei geradem Wert, wenn der zur Wurzel $u=0$ gehörende Koeffizient ausgeschlossen wird, die Form
\[
\gamma u^{n-2}+\gamma'u^{n-4}+\gamma''u^{n-6}+\mathrm{usw.}
\]
hat. Das Beispiel des Artikels 12 zeigt genau diese Reduktion für $n=6$. Entgegengesetzten Werten von $u$ entsprechen somit offenbar stets gleiche Koeffizienten. Im Fall, dass $n$ gerade ist, lässt sich außerdem der zur Wurzel $u=0$ gehörende Koeffizient leicht allgemein a priori angeben. Man erhält diesen Koeffizienten, indem man in $\frac{U'}{(\frac{\mathrm{d}U}{\mathrm{d}u})}$ den Wert $u=0$ einsetzt. Den Wert des Zählers $U'$ für $u=0$ haben wir bereits in Artikel 18 angegeben; der Wert des Nenners ergibt sich aber aus demselben Artikel zu
\[
\begin{aligned}&=\pm\frac{3.5.7\mathinner{.\,.\,.\,.\,.}(n+1)}
{(n+3)(n+5)\mathinner{.\,.\,.\,.\,.}(2n+1)}\\[5pt]&=\pm\frac{3.3.5.5.7.7\mathinner{.\,.\,.\,.\,.}(n+1)(n+1)}
{3.5.7.9.11\mathinner{.\,.\,.\,.\,.}(2n+1)}\end{aligned}
\]
und somit lautet der gesuchte Koeffizient
\[
=\left(\frac{2.4.6.8\mathinner{.\,.\,.\,.\,.} n}
{3.5.7.9\mathinner{.\,.\,.\,.\,.}(n+1)}\right)^2
\]

\SeventhArticle{21}
Das Polynom in $u$, welches die Koeffizienten $R,R',R''$ usw. in dem hier betrachteten Fall darstellt, lässt sich auch unabhängig vom allgemeinen Verfahren des Artikels 11 folgendermaßen finden. Wir differenzieren zunächst die Gleichung
\[
\varphi-\frac{U'}U=\frac{U''}U
\]
Nach anschließender Substitution von $\frac{\mathrm{d}\varphi}{\mathrm{d}u}
=\frac1{1-uu}$ und Multiplikation mit $UU(uu-1)$ erhalten wir
\[
\begin{aligned}&(uu-1)U'\frac{\mathrm{d}U}{\mathrm{d}u}
-U\left(\frac{\mathrm{d}U'}{\mathrm{d}u}.(uu-1)+U\right)\\[5pt]&=(uu-1)UU\frac{\mathrm{d}(\frac{U''}U)}{\mathrm{d}u}\end{aligned}
\]
\SourceLeaf{201}
Die Glieder auf der linken Seite dieser Gleichung bilden offenbar ein Polynom in $u$; daher müssen sich auf der rechten Seite die Koeffizienten der Potenzen von $u$ mit negativen Exponenten notwendig gegenseitig aufheben.

Doch $\frac{\mathrm{d}\frac{U''}U}{\mathrm{d}u}$ ergibt eine unendliche Reihe, deren erstes Glied lautet
\[
-\left(\frac{1.2.3.4\mathinner{.\,.\,.\,.\,.}(n+1)}
{1.3.5.7\mathinner{.\,.\,.\,.\,.}(2n+1)}\right)^2u^{-(2n+4)}
\]
Wird diese Reihe also mit $(uu-1)UU$ multipliziert, so kann sich nichts anderes als die konstante Größe
\[
-\left(\frac{1.2.3.4\mathinner{.\,.\,.\,.\,.}(n+1)}
{1.3.5.7\mathinner{.\,.\,.\,.\,.}(2n+1)}\right)^2
\]
ergeben. Daraus schließen wir\HistoricalNote{20}, dass
\[
(uu-1)U'\frac{\mathrm{d}U}{\mathrm{d}u}
+\left(\frac{1.2.3.4\mathinner{.\,.\,.\,.\,.}(n+1)}
{1.3.5.7\mathinner{.\,.\,.\,.\,.}(2n+1)}\right)^2
\]
durch $U$ teilbar ist; daher ist dem rationalen Ausdruck $\frac{U'}{(\frac{\mathrm{d}U}{\mathrm{d}u})}$, der die Koeffizienten $R,R',R''$ usw. liefert, das folgende Polynom äquivalent:
\[
-\left(\frac{1.3.5.7\mathinner{.\,.\,.\,.\,.}(2n+1)}
{1.2.3.4\mathinner{.\,.\,.\,.\,.}(n+1)}U'\right)^2.(uu-1)
\]
Anstelle dieser Funktion, die vom Grad $2n+2$ ist und offenbar nur gerade Potenzen von $u$ enthält, kann der bei ihrer Division durch $U$ entstehende Rest genommen werden; er ist vom Grad $n$ beziehungsweise $n-1$, je nachdem, ob $n$ gerade oder ungerade ist. Wollen wir jedoch im ersteren Fall den zur Wurzel $u=0$ gehörenden Koeffizienten lieber ausschließen, so nehmen wir anstelle jener Funktion ihren Rest bei der Division durch $\frac Uu$, der nur bis zum Grad $n-2$ reicht.

